\section{Limitations and threats to validity}
\label{sec:limitations}

This revision establishes an implemented governance path, a local hostile falsifier suite, a frozen prospective design, and one diagnostic attribution archive. It does not establish transferable protected improvement, reduced harmful transfer, or an integrity advantage against matched self-edit or self-evolution baselines.

\paragraph{The theory is wider than the evidence.}
The revision-factorization, Bayes-risk and discriminator-cover results apply to
arbitrary finite decision alphabets and, for the cover theorem, finite latent
fibres with deterministic protected test outcomes. The stochastic
transcript-law theorem characterizes a fixed adaptive policy. The
state-indexed frontier additionally enumerates optimal finite-horizon policies,
but only for finite observable Markov-sufficient laboratory state, finite legal
actions, known kernels and declared losses. It does not validate those kernels,
identify unknown kernels from data, solve the unrestricted infinite-horizon
problem, or supply an unknown-kernel sample-complexity bound. The results do not
show that the eight revision classes are complete,
that the true latent situation is in the registered fibre, or that the required
protected distributions are identifiable from finite samples. A
registered cover proves structural separability under its model, not empirical
power, calibration, transfer or beneficial revision.

\paragraph{What is not claimed as novel.}
This is not the first self-improving agent, the first evolutionary program-search system, the first candidate archive, the first issue-centric self-improvement state, or the first failure-driven adaptation loop. Meta-agent design search, archive-based self-modification, persistent issue state, multi-hypothesis attribution, counterfactual repair, source-level self-rewriting with replay validation, pre-commit gating, evaluator locking and anytime-valid acceptance are all prior work; they are discussed and cited in Section~\ref{sec:related}, and each is absorbed rather than claimed. What is left as a candidate difference is the composition: the ordered, non-compensatory sequence from persistent issue state through causal discrimination, invention readiness, isolated execution, replay, independent protected fresh transfer and protected assurance, with promotion authority structurally outside the evaluated system. That is a claim about the safety and reliability of a development process, and architecture alone does not establish it.

\paragraph{The diagnostic archive does not test the primary hypothesis.}
The archive scores 21/24 on hidden-cause labels. That quantity is a
single-model, single-run descriptive accuracy. Its nominal Wilson interval
treats 24 fixed constructed cases as Bernoulli units; without a probability
sampling frame or a licensed dependence model it is not a population confidence
interval. It does not measure protected fresh-task improvement, harmful
transfer, or false protected acceptance. The three residual errors remain
errors.

\paragraph{The 24-case suite is not an execution-frozen campaign.}
The suite covers eight families, but evaluator hashes and several fresh-task hashes remain placeholders. Motivating, replay and fresh split identities, provider and model revisions, and the protected evaluator epoch are all unbound in both frozen protocol revisions.

\paragraph{The wide source universe is not bound.}
An outcome-blind successor fixes 16 candidate roots across two waves, with
four metadata providers, four artifact modalities, four nominal domains and
two roots per domain in each wave. Its 13/13 conformance checks and 16/16 live
identity checks close metadata structure only. Eight roots lack exact
historical Crossref bytes; content-class rights and case eligibility are
unresolved for all sixteen; and no project/issue case, protected revision
label or outcome has been opened. The 768-cluster successor therefore remains
unexecuted.

Public bug corpora widen development coverage but do not create protected
evidence. SWE-bench Multilingual, Defects4J, BEARS, GitBug-Java, Bugs.jar,
BugSwarm, BugsInPy, ManyBugs, QuixBugs and SWE-bench Verified now have metadata
identities, but framework or repository licences do not automatically license
every upstream project, issue, patch, test, log, dependency or container byte.
None supplies protected eight-class revision responsibility or protected fresh
transfer. Payload access therefore remains conditional on a per-project and
per-content-class rights/eligibility census performed without outcome-based
source selection.

\paragraph{No matched baselines or ablations were executed.}
The six comparator identities are source-bound, but no matched comparison has
been executed. SWE-agent, MOSS, Darwin G\"odel Machine, ADIAS, Double Ratchet's
metric-only arm and ScienceClaw each have an exact paper/repository revision
and native entrypoint. A generic refined adapter contract is now total on its
declared synthetic domain. C1 SWE-agent now has an outcome-blind native
parser/write-isolation and partial resource binding but still has 12 blocking
fields. C2 MOSS separately has a native parser, closed fallback policy,
wallclock and dependency locks but retains 14 blockers; its authoritative
source tree omits the benchmark runtime required by its native trial runner.
C3 DGM has a closed patch parser, fallback policy and outer wallclock but
retains 15 blockers, including a resolved dependency lock: its 23 declarations
have zero exact pins. C4 ADIAS has a closed report parser, fallback/search
policy and outer wallclock but retains 15 blockers, including dependency/image
identity, task rights and rootless isolation. C5 Double Ratchet has an isolated
write surface, parser, resource envelope and dependency lock but retains 12
blockers because the official loop regenerates solver outputs and reuses its
development-locked metric every round. C6 ScienceClaw has a strict draft-shape
parser and outer wallclock but retains 16 blockers: no native revision selector,
no close-all fallback switch, no locked 112-declaration/lazy-dependency
environment, and no complete resource, rights or custody binding. None has a
protected panel, complete matched resource manifest or result-bearing archive
on this subject.
Identity closure therefore narrows the blocker from which systems to whether
their native interfaces can be translated lawfully and executed under equal
exposure without protected feedback. It supplies no effect estimate, and the
matched baseline and ablation comparison remains undetermined rather than
null.

\paragraph{Adapter conformance is not comparator performance.}
The six-arm V2 adapter audit proves a formal interface boundary and passes
90/90 synthetic contract cases; V3 extends conformance to 231/231 declared
fictional cases. Neither uses a native output example or outcome. All six raw
same-visible-symptom fibres are mixed, so zero raw native statuses license a
singleton class. V3's 40 singleton case records occupy 20 supported constant
fibres only after a host-validated class certificate and exclusive-front write
proof are supplied; the other 191 cases remain \texttt{UNRESOLVED}. These
records show only that a proposed action can be typed when the missing
identifying information is supplied. They do not show that the action is
correct, minimal, preserving, transferable, or safe. ScienceClaw has no native
singleton support, Double Ratchet conditionally supports only evaluator repair,
and the V3 ledger's 108 blocking fields are preserved as its historical result.
The C1 and C5 V4 packets each resolve six fields, while C2, C3 and C4 resolve
respectively four, three and three fields, while C6 resolves two. This leaves
12 each for C1/C5 plus 14 for C2 plus 15 each for C3/C4 plus 16 for C6, or 84
current blocker instances.
Zero of six arms has a complete
matched resource, rights, arm-native adapter and protected-scorer manifest.

\paragraph{Root-cause compression is not blocker closure.}
The panel-wide V5 registry proves that the 84 current blocker instances are an
exact partition into 14 equivalence classes and five evidence-producing work
packages. This permits shared acquisition where the same evidence is logically
sufficient. V5 itself does not permit a schema to fan out as evidence and its
arm-field delta is zero. V6 subsequently provides one substantive
rights-cleared visible-case packet and exact per-arm acceptance receipts,
closing its targeted 12 fields and reducing the census to 72 blockers. It does
not by itself close the six native task environments or the 18-field
independent-custody root. V7 later closes only C1's environment with actual
offline setup/configuration bytes, reducing the census to 71 blockers and the
environment root to five arm-specific bundles. Moreover, the wallclock
bindings are not matched: C1's 3,600-second run
and 30-second grace differ from C2--C6's 21,600-second run and 120-second grace,
and C5 has an extra 10-second candidate subprocess limit. Equal nominal field
status is therefore insufficient for equal effective exposure.

The C3 V8 candidate-safe seed does not close its environment. Its unchanged
native parent-selection path is outcome-indexed: a fresh singleton
\texttt{initial} parent is skipped unless prior performance and identifier
partitions exist. V9 establishes this boundary only for the exact released DGM
commit. It does not prove that outcome-free Darwin G\"odel Machine
initialization is impossible in future designs. A future source-native identity
element can reopen the gate; a locally invented score-free selector is instead
a new comparator requiring separate preregistration and matched-panel review.

\paragraph{Fresh source-packet assurance is not uniform.}
The C4 validator includes the hash of metadata containing a random temporary
generation path in its audit receipt. Re-running it changes
\texttt{AUDIT\_RECEIPT\_V4.json}, so the frozen C4 packet checksum does not
reproduce without a path-invariant repair. The C6 validator requires a cleaned
\texttt{.source-audit} checkout that is no longer present, so its fresh
source-level rerun is \textsc{CannotCheck}. The field-registry, result and
parser bytes used for the 42/84 census remain directly hash-verifiable; that
narrower assurance must not be expanded into full packet reproducibility.

\paragraph{Common content rights are not native-environment rights.}
The V6 Lang-1 core is substantive and lawful for its declared candidate-visible
task. Its exact acceptance is conditional on every adapter using only that
core. C1's image/setup/run configuration is now bound. C2's benchmark companion
and native session, C4's domain/tasks/seeds/limits, C5's solver
outputs/prompts/membership and C6's profile/skills/tools/fallback closure remain
absent. C3 additionally lacks an outcome-free native initial parent; its
host/agent split and endpoint bytes alone cannot supply one. Reintroducing
bundled native content triggers a new
component-rights audit rather than inheriting the common grant.

\paragraph{The released MOSS runtime is incomplete at the pinned identity.}
The paper-linked commit, tree and release archive are byte-addressed, and the
root/OpenClaw licence layers are retained. However, the native trial runner
imports and writes under \texttt{benchmark/claw-eval}, while the complete
non-truncated authoritative tree contains no \texttt{benchmark/} path and its
setup instructions do not fetch one. Source availability therefore does not
establish task-runtime or benchmark rights. This absence is retained as
\textsc{CannotCheck}; it is not repaired with invented fixtures, unversioned
bytes or a different evaluator.

\paragraph{The C2 successors are lawful components, not an assembled MOSS arm.}
V11, V12 and V13 close different fields for three distinct authored successor
identities. V11's runtime and V12's source-core rights are not inherited by
V13; V11 likewise does not inherit V12 or V13, and V12 does not inherit V11 or
V13. Adding 8/21, 9/21 and 8/21 states would therefore manufacture an
unexecuted fourth identity. V13's scratch probe establishes a complete
nine-file image rights/SBOM chain and a constrained container runtime receipt,
but that probe is not the missing released MOSS benchmark runtime and does not
bind \texttt{runtime.container\_or\_environment} for released MOSS. Released
MOSS remains 7/21 bound and 14 blocking, the 55/126 panel census is unchanged,
and no performance, H1--H4 or superiority inference follows.

\paragraph{The released DGM runtime is not a frozen execution environment.}
The source commit, tree, uncompressed Git-archive digest and Apache-2.0 root
licence are byte-addressed. That does not bind execution: all 23 Python
dependency declarations are unpinned, the Docker base image is mutable, and
the candidate source includes outcome/initial paths that must be excluded
before access. The released outer-loop timeout is called only after a future
has completed, and the source retains an argparse choice concatenation defect
and a destructive-untrusted-code warning. These findings are source preflight,
not a failed DGM performance observation; native execution remains
\textsc{CannotCheck}.

\paragraph{The released ADIAS runtime is not a rights-closed isolated execution environment.}
The source commit, tree, prefixed Git-archive digest, licence and 1,578-file
manifest are byte-addressed. The root source licence is CC-BY-NC-SA-4.0 and
does not itself close the exact intended-use rights; bundled task data expose
no component licence/NOTICE files under \texttt{data/}. Four VCS declarations
are unpinned, the base image is mutable, native containers use host networking
with provider keys, and a root agent defeats the source's \texttt{chmod 0}
domain-hiding convention as a security boundary. Same-run final testing is not
an external protected scorer. These are source-preflight defects, not failed
ADIAS performance observations; native execution remains
\textsc{CannotCheck}.

\paragraph{The released Double Ratchet metric loop is not a protected scorer.}
The source, entrypoint, licence, 113-blob tree, prefixed Git archive and
46-entry dependency lock are byte-addressed, and V4 adds a results-only write
surface and resource envelope. Yet the official loop regenerates hosted solver
outputs and evaluates, persists and prints \texttt{eval\_locked} each round.
The tree contains no frozen split/result payload, and its MBPP/report-generation
surface is not a P5 dossier/eight-class task adapter. A development-only
surrogate may remain inside the loop, but the final fresh panel must be external
and one-shot. These are interface/custody findings, not a failed Double Ratchet
performance observation; native P5 execution remains \textsc{CannotCheck}.

\paragraph{The residual-error successors retain distinct execution status.}
The three factorial discriminators were derived from retained errors and are
therefore new studies, not confirmatory rescoring of the original 24 cases.
P5-RD-01 and P5-RD-03 have no observations. P5-RD-02 V1 remains
\textsf{CANNOT\_CHECK} with an adverse-cell modifier; V2 records one
preregistered public-development known-fix cluster with a within-cluster
\textsf{IMPLEMENTATION\_MAIN\_EFFECT}; and V3 is a post-outcome audit-only
archival replay. V2 does not rewrite V1, V3 does not recover V2's overwritten
checkout streams, and none changes P5-HC-012 or supplies protected freshness.
A main effect, interaction, joint null, adverse terminal, or persistent
observational equivalence remains attached to its own study identity; none may
be collapsed into a post-hoc success.

\paragraph{The 24/24 V2 attribution result is an instrument-stage diagnosis.}
The V2 study reuses the same public 24-case suite and changes the attribution
instrument from free-form causal re-derivation to staged evidence extraction
plus deterministic diagnosis. Its control reproduces 21/24 and its treatment
reaches 24/24, but protocol and outcomes first appear in the same commit;
independent preregistration chronology is therefore \textsc{CannotCheck}.
The requested GLM-5.2 identity also differs from the recorded GLM-5.3 served
identity. Without signed provider custody, a blind fresh suite, source-disjoint
transfer or a matched protected panel, this result cannot establish model
improvement, general attribution accuracy, H1--H4 or superiority.

\paragraph{The frozen revision-level panel exposes an identification mismatch.}
All bound execution identities verified and the independent scorers agreed,
but FULL\_T7 observed one discriminator while each hypothesis expectation map
required two. The resulting $0.125$ accuracy and fresh-transfer rate are
therefore a measured property of the frozen one-probe mechanism on this panel,
not a software failure and not a result that may be erased by changing the
panel after outcome access. The generic-diagnosis and random controls outperform
the subject on this constructed panel. A successor may test completable
expectation sets and explicit preservation conflicts, but it must be separately
frozen; it cannot retroactively turn this no-terminal result into support for
H1 or general superiority. That successor was separately frozen and executed;
its result is reported in Section~\ref{sec:results} and it does not alter this
paragraph: the $0.125$ no-terminal outcome on the V3 panel stands, and the
successor's own fresh-transfer measure is noninferior to its strongest
comparator rather than superior, so H1 and general superiority remain open.

\paragraph{No live development cycle is bound.}
The live-trial packet still records an unbound corpus revision. Provider and protected-verifier credentials were unset for this revision, so no live campaign was started. A previously advertised report of a flawless run is not used: its artifact identity and repository identity diverged before integration.

\paragraph{Statistical threats.}
The diagnostic contains 24 fixed constructed cases from one model run, not 24
independent campaigns or a probability sample. P5-RD-02 contains one
source/bug/fix cluster; its environments, cells and phase streams are repeated
within-cluster measurements, not independent replicates. Neither archive can
support the predeclared five-percentage-point fresh-improvement target or the
two-percentage-point harm guard. There is no paired cluster-level bootstrap
against a baseline, no multiplicity-adjusted secondary family, and no tail or
family regression from a campaign. The macro-F1 reported in the original
sidecar assumed flawless precision and is not promoted.

\paragraph{Contamination and evaluator custody.}
The diagnostic run used locally visible gold labels to score attributions after the fact. That is admissible for an offline label-recovery diagnostic. It is not admissible as protected-evaluator evidence: the candidate-facing study still requires labels, fresh payloads and scoring rubrics to remain outside challenger custody.
